% !TEX root = ../thesis.tex

\chapter{Návrh riešenia}
\label{ch:navrh-riesenia}

V tejto kapitole je opísaný návrh pipeline doménovej adaptácie malých jazykových modelov pre energetické úlohy: od výberu modelov a prípravy dát až po metodologické východiská prístupu QLoRA a návrh inferenčnej a vyhodnocovacej vrstvy.

\section{Výber modelov a ich architektonické vlastnosti}

Pre experimenty boli zvolené tri otvorené modely triedy SLM (Small Language Model), ktoré sú veľkosťou porovnateľné, ale líšia sa architektonickými detailmi:

\begin{itemize}
    \item \textbf{Phi-3 Mini (3.8B)} (\texttt{microsoft/Phi-3-mini-4k-instruct}),
    \item \textbf{Gemma 3 4B} (\texttt{google/gemma-3-4b-it}),
    \item \textbf{Qwen3 4B} (\texttt{Qwen/Qwen3-4B-Instruct-2507}).
\end{itemize}

Hlavné kritériá výberu:
\begin{itemize}
    \item veľkosť približne 4B parametrov, ktorá umožňuje tréning na obmedzených GPU zdrojoch pri 4-bitovej kvantizácii,
    \item dostupnosť oficiálnych instruction/chat verzií,
    \item kompatibilita s ekosystémom \texttt{transformers + peft + trl + bitsandbytes},
    \item architektonická rôznorodosť vhodná na porovnanie (odlišná tokenizácia, chat šablóny, mechanizmy attention a implementácia forward priechodu).
\end{itemize}

Z inžinierskeho pohľadu sa modely vhodne dopĺňajú: Phi-3 Mini je orientovaná na vysokú efektivitu pri malom počte parametrov, Gemma 3 4B dosahuje kvalitné sledovanie inštrukcií a pri textovom tréningu vyžaduje špecifickú integráciu (keďže ide o multimodálnu architektúru), Qwen3 4B ponúka stabilný výkon v analytických úlohách a konverzačnom režime.

\subsection{Prečo Instruct verzie namiesto Base váh}

V projekte je cieľovým scenárom \textit{interaktívny asistent}, kde používateľ zadáva inštrukciu a model má vrátiť stručnú, štruktúrovanú a prakticky použiteľnú odpoveď. Pre takýto režim sú Instruct váhy vhodnejšou inicializáciou než Base váhy z týchto dôvodov:

\begin{itemize}
    \item Instruct modely už prešli učením vo formáte \enquote{inštrukcia $\rightarrow$ odpoveď}, preto sa v SFT zvyčajne učia rýchlejšie a stabilnejšie,
    \item Base modely sú optimalizované hlavne na pokračovanie textu, nie na disciplinované konverzačné odpovede,
    \item pri obmedzenom výpočtovom rozpočte je výhodné nepreúčať základné správanie modelu od nuly, ale sústrediť sa na doménové znalosti,
    \item štart z Instruct váh znižuje riziko degradácie konverzačných schopností počas PEFT tréningu na menšom doménovom datasete.
\end{itemize}

Inými slovami, Instruct inicializácia v tomto projekte funguje ako silný \textit{behavioral prior}: QLoRA adaptéry dolaďujú doménové znalosti a štýl odpovedí, nie základnú schopnosť modelu nasledovať inštrukcie.

\section{Dáta a štruktúra datasetu}

\subsection{Formát a organizácia dát}

Tréning bol realizovaný na jednom JSONL datasete v štruktúrovanom formáte:
\begin{itemize}
    \item \texttt{data/energy\_data\_structured.jsonl} -- finálny tréningový dataset pre cieľový formát odpovede.
\end{itemize}

Každý riadok JSONL predstavuje jednu tréningovú vzorku s poľami:
\begin{itemize}
    \item \texttt{prompt}: otázka/inštrukcia používateľa,
    \item \texttt{completion}: cieľová odpoveď asistenta.
\end{itemize}

Formálne je štruktúra vzorky definovaná ako
\begin{equation}
    x_i = \{p_i, c_i\}, \quad p_i \in \mathcal{P},\; c_i \in \mathcal{C}
\end{equation}
kde $p_i$ je používateľský vstup a $c_i$ je cieľová odpoveď v doménovom štýle.

V použitej verzii projektu má dataset nasledujúci rozsah:
\begin{itemize}
    \item \texttt{energy\_data\_structured.jsonl}: 1159 záznamov.
\end{itemize}

Pre preferenčnú optimalizáciu je navyše pripravený samostatný dataset:
\begin{itemize}
    \item \texttt{energy\_data\_dpo.jsonl}: 1126 záznamov vo formáte \texttt{prompt/chosen/rejected}.
\end{itemize}

\subsection{Logika prompt-completion a jej vhodnosť}

Zvolený formát prompt-completion je pre tento scenár vhodný z troch dôvodov:

\begin{itemize}
    \item \textbf{Priama kompatibilita s Causal LM.} Dialóg sa serializuje cez chat-template tokenizéra do sekvencie system/user/assistant a model sa učí predikovať tokeny asistenta.
    \item \textbf{Kontrolovaný tréningový signál.} V pipeline sa používajú maskované labely: tokeny systémovej a používateľskej časti majú hodnotu $-100$, takže loss sa počíta iba z completion časti. Tým sa model netrestá za reprodukciu vstupu a sústredí sa na kvalitu odpovede.
    \item \textbf{Stabilita pre PEFT.} Pri QLoRA sa aktualizuje len obmedzený počet parametrov (LoRA adaptéry), preto je dôležitý jednoznačný supervised signál. Štruktúrovaný completion tento signál poskytuje.
\end{itemize}

V štruktúrovanej verzii odpovede sa používa šablóna:
\begin{verbatim}
ANSWER: ...

KEY FACTS:
- ...
- ...

RISK LEVEL: Low/Medium/High
...

CONFIDENCE: High/Medium/Low
...
\end{verbatim}

Táto štruktúra je prínosná nielen pre používateľa, ale aj pre samotné učenie: model si osvojuje opakovateľný formát generovania, čo zvyšuje predvídateľnosť výstupu počas inferencie.

\subsection{Príprava dát v zdieľanom dátovom loaderi}

Implementácia SFT používa spoločnú dátovú vrstvu pre všetky tri modely (\texttt{training/shared/data\_loader.py}). Táto vrstva zabezpečuje:
\begin{itemize}
    \item načítanie JSONL datasetu cez knižnicu \texttt{datasets},
    \item serializáciu dialógu \texttt{system/user/assistant} cez \texttt{tokenizer.apply\_chat\_template},
    \item vytvorenie \texttt{input\_ids}, \texttt{attention\_mask} a \texttt{labels},
    \item maskovanie všetkých netréningových tokenov hodnotou $-100$,
    \item rozdelenie datasetu na trénovaciu a validačnú časť.
\end{itemize}

Výhodou tohto riešenia je jednotná metodika spracovania dát naprieč modelmi. Rozdiely medzi experimentmi sú preto dané najmä architektúrou modelov a hyperparametrami, nie odlišným predspracovaním.

\section{Metodológia QLoRA}

\subsection{Technická schéma metódy}

QLoRA kombinuje dva mechanizmy:
\begin{itemize}
    \item 4-bitovú kvantizáciu základného modelu (v projekte: NF4 cez \texttt{BitsAndBytesConfig}),
    \item low-rank adaptáciu pomocou LoRA matíc trénovaných nad zmrazenými kvantovanými váhami.
\end{itemize}

Nech má lineárna vrstva váhovú maticu $W$. V LoRA je aktualizácia definovaná ako
\begin{equation}
    W' = W + \Delta W, \qquad \Delta W = BA,
\end{equation}
kde $A \in \mathbb{R}^{r \times d}$, $B \in \mathbb{R}^{k \times r}$ a $r$ je malý rank ($r \ll \min(d,k)$).

V QLoRA sú pôvodné váhy $W$ uložené v 4-bitovej reprezentácii, no výpočty a učenie adaptérov prebiehajú vo vyššej presnosti (v projekte prevažne \texttt{bfloat16}). Tým sa znižujú nároky na VRAM pri zachovaní numerickej stability potrebnej na doučenie.

\subsection{Prečo je QLoRA efektívna pri obmedzených zdrojoch}

Efektivita vyplýva z kombinácie viacerých faktorov:
\begin{itemize}
    \item \textbf{Nižšie pamäťové nároky na váhy:} prechod zo 16-bitovej na 4-bitovú reprezentáciu zmenšuje pamäťový footprint modelu približne 4-násobne.
    \item \textbf{Malý počet trénovateľných parametrov:} učia sa len LoRA adaptéry, nie celé váhy modelu.
    \item \textbf{Optimalizátor pre low-precision režim:} použitý \texttt{paged\_adamw\_8bit} znižuje nároky na pamäť stavov optimalizátora.
    \item \textbf{Gradient accumulation:} umožňuje väčší efektívny batch bez nárastu okamžitej spotreby VRAM.
\end{itemize}

V praxi to umožňuje doučenie modelov triedy 4B na jednom GPU spotrebiteľskej triedy, čo zodpovedá hardvérovým obmedzeniam projektu.

\section{Implementácia inferencie a aplikačnej vrstvy}

Okrem tréningovej pipeline bola v projekte implementovaná aj produkčne orientovaná inferenčná vrstva postavená na \texttt{FastAPI}. Aplikácia je organizovaná modulárne (\texttt{core}, \texttt{api}, \texttt{services}, \texttt{models}, \texttt{templates}, \texttt{static}), čo zjednodušuje rozšíriteľnosť aj údržbu.

Kľúčovou súčasťou je služba \texttt{ModelService}, ktorá:
\begin{itemize}
    \item dynamicky načítava zvolený model (Phi-3, Gemma 3, Qwen3),
    \item podľa konfigurácie pripája príslušný LoRA adaptér,
    \item pri prepnutí modelu uvoľní predchádzajúci model z VRAM,
    \item štandardizuje generovanie odpovedí cez jednotné parametre dekódovania.
\end{itemize}

Pre stabilné ukončenie generovania sa používajú modelovo špecifické stop tokeny (napr. \texttt{<|end|>} pre Phi-3, \texttt{<|im\_end|>} pre Qwen, \texttt{<end\_of\_turn>} pre Gemma). Výstup je následne čistený od technických tokenov a pri Qwen modeli aj od prípadných \texttt{<think>} blokov.

API vrstva poskytuje najmä endpointy \texttt{/api/chat}, \texttt{/api/models} a \texttt{/api/health}. Komunikácia prebieha vo formáte chat správ (história + nový prompt), pričom systémový prompt je centrálne definovaný v konfigurácii a je jednotný pre všetky modely. Tým sa zabezpečuje porovnateľnosť odpovedí medzi modelmi aj konzistentné používateľské rozhranie.

\section{Vyhodnocovacia pipeline a reprodukovateľnosť}

Pre experimentálne porovnanie modelov je pripravený samostatný skript \texttt{scripts/evaluate\_models.py}, ktorý meria kvalitatívne aj prevádzkové metriky:
\begin{itemize}
    \item G-Eval (relevance, coherence, fluency, factual accuracy),
    \item perplexitu na testovacích vzorkách,
    \item rýchlostné metriky \texttt{TTFT} a \texttt{TPS},
    \item spotrebu VRAM pri rôznych dĺžkach kontextu.
\end{itemize}

Reprodukovateľnosť tréningu je podporená oddelením konfigurácií podľa modelových rodín (\texttt{training/phi3}, \texttt{training/gemma3}, \texttt{training/qwen3}) a jednotnými vstupnými dátami. Všetky experimenty ukladajú výsledné adaptéry do priečinkov \texttt{final\_adapter}, čo umožňuje ich okamžité použitie v API vrstve bez dodatočných konverzií.

\section{Zhrnutie návrhu riešenia}

Návrh riešenia je koncipovaný ako reprodukovateľná a výpočtovo efektívna pipeline: Instruct modely poskytujú vhodné počiatočné správanie asistenta, zdieľaný dátový loader zabezpečuje jednotný supervised signál, QLoRA umožňuje doménovú adaptáciu na obmedzenom hardvéri a DPO dopĺňa preferenčnú kalibráciu odpovedí. Výsledkom je praktický kompromis medzi kvalitou odpovedí, výpočtovou náročnosťou a rýchlosťou experimentálnych iterácií spolu s priamo nasaditeľným API riešením.
